\subsection{Batching a large-rank projector edge}
\label{sec:projector-batching}

The usual matrix-shear compiler uses one smaller interchange per rank
unit. For a projector of rank close to the ambient dimension, most of
these pivot interchanges can instead be one interchange of two
contiguous longer fields. The construction below does not gather
separated fields or reverse their digits.

\paragraph{A simultaneous rational change of basis.}
Let $\mathcal P$ be a fixed finite family of rational idempotents of
dimension $m$. For $P\in\mathcal P$ with
$m/2<a:=\operatorname{rank}P<m$, put $r=m-a$. There is one
$S\in\operatorname{GL}_m(\mathbb Q)$ such that, for every such $P$,
the submatrix of $SPS^{-1}$ on its first $r$ rows and last $r$ columns
is invertible.

To prove this, first consider one rank-$a$ idempotent. Every rational
idempotent is rationally similar to $\operatorname{diag}(I_a,0)$,
by concatenating rational bases of its image and kernel. Put
$t=m-2r>0$. On the first and last $r$ coordinates take $r$ disjoint
rank-one blocks
\[
 \frac12\begin{pmatrix}1&1\\1&1\end{pmatrix},
\]
pairing coordinate $i$ with $m+1-i$, and take the identity on the middle
$t$ coordinates. This is a rational rank-$a$ idempotent whose indicated
corner is invertible. Hence the determinant of that corner of $SPS^{-1}$
is not the zero rational function of the entries of $S$. Clear powers
of $\det S$ from its denominator. The product of the resulting nonzero
polynomials for the finite family, and $\det S$, is nonzero over
$\mathbb Q$. A nonzero polynomial over the infinite field $\mathbb Q$
cannot vanish at every rational point. This proves simultaneous existence.

The choice is effective: enumerate integer matrices and retain the first
invertible one satisfying all the indicated rational determinant tests.
This search terminates. The family and $m$ are fixed, so this finite
calculation is part of the fixed-machine preparation. Conjugate every
terminal and gate matrix by this same $S$. Common gate matrices,
every edge rank, the scalar circuit, and the endpoint identity are
preserved, since $SIS^{-1}=I$. The address prime is then chosen after
this conjugation and the triangular factorizations, avoiding their
finitely many denominators and invertible diagonal numerators as in the
original construction.

Only the edges selected for batching need belong to $\mathcal P$.
All other matrices can continue to use the original rank-by-rank
compiler. In the present finite network, every selected edge is a
nonzero scalar multiple of an idempotent: for nested nondegenerate
labels $U\subset V$,
\[
 (P_V-P_U)^2=P_V-P_U,
\]
and reversing the edge only changes its sign. Scalar multiples do not
change the pivot locations and are absorbed into the triangular factors.

\paragraph{The contiguous middle pivots.}
Let $P^2=P$ have rank $a>m/2$, let $r=m-a>0$, and assume its
first-$r$-row, last-$r$-column corner is invertible. Split the ordered
coordinates into
\[
 L=[1,r],\qquad M=[r+1,m-r],\qquad R=[m-r+1,m].
\]
The complement $Q=I-P$ has rank $r$. Write $Q=UV$, where $U$ has
$r$ columns, $V$ has $r$ rows, and $VU=I_r$. Such a factorization
comes from a basis of $\operatorname{im}Q$. As $L$ and $R$ are disjoint,
\[
 P_{L,R}=-U_LV_R.
\]
Its invertibility makes both square matrices $U_L,V_R$ invertible.

Run the lower-lower elimination in the retained physical order. Its
first $r$ pivots occupy the rows $L$ and the columns $R$, in some
permutation. Indeed their corner is invertible, so at each of these
steps the next row has a nonzero entry in an as-yet active column of
$R$, which precedes all other candidates in the rightmost-pivot rule.
The active Schur complement on rows $M\cup R$ and columns $L\cup M$
has middle block
\begin{align*}
 P_{M,M}-P_{M,R}P_{L,R}^{-1}P_{L,M}
 &=I-U_MV_M
   -(-U_MV_R)(-V_R^{-1}U_L^{-1})(-U_LV_M)\\
 &=I_{m-2r}.
\end{align*}
Thus the next $m-2r$ pivots are precisely the diagonal positions
$(i,i)$ with $i\in M$. Columns $L$ lie to their left and are cleared
by the allowed lower-triangular column operations. The remaining
block has rank zero, because the number of isolated pivots is already
$r+(m-2r)=m-r=a$.

Consequently the factorization $P=E_1\Pi E_2$, with $E_1,E_2$
invertible lower triangular, has exactly $r$ corner pivots and one
contiguous identity block of size
\[
 t=m-2r=2a-m.
\]
The same conclusion holds for $cP$ with any fixed nonzero rational $c$.
No assertion that an arbitrary rank-$a$ matrix has this pivot profile
is needed; the idempotent identity is what makes the middle block
exactly an identity matrix.

\paragraph{One call for the identity block.}
Subdivide each original contiguous chunk into $m$ consecutive fields
of width $b$. The triangular factors still cost $O(V)$ in the retained
later-field affine primitive. Each of the $r$ corner pivots uses one
ordinary width-$b$ interchange. For all $i\in M$ together, perform
\[
 D_i\leftarrow D_i+H_i,
 \qquad
 (H_M,D_M)\leftarrow(D_M,H_M),
 \qquad
 D_i\leftarrow H_i-D_i,
\]
over $\mathbb Z/q^b\mathbb Z$.
The first and last maps are coordinatewise later-field updates, each
using a fixed number of linear scans. The middle map is a single
interchange of the two contiguous width-$tb$ chunks $H_M,D_M$.
It preserves the order of their $b$-digit fields and hence implements
all diagonal pivots exactly. The intervening fields are spectators.
In particular, addition is componentwise modulo $q^b$, not addition
with carries across the concatenated width-$tb$ chunks.

The edge cost therefore has $r$ recursive calls of width $b$ and one
of width $tb$, in place of $a$ calls of width $b$, with the same
$O(V)$ nonrecursive work. If $a=m$, one may leave the edge in its
original $m$-singleton form. This avoids any same-width recursion.
For every batched edge in this paragraph, $0<t\le m-2$, so every
child contracts its width by a fixed factor less than one.

\paragraph{The resulting characteristic expression.}
Let $\mathcal B$ be the selected large-rank projector edges, let
$r_e=m-a_e$ and $t_e=m-2r_e$ for $e\in\mathcal B$, and retain the
total rank sum $s$. For an exponent $0<\tau<1$, the volume-weighted
recursive contribution is
\begin{equation}\label{eq:batched-characteristic}
 \mathcal C(\tau)
 =\frac1W\left[
   \left(s-\sum_{e\in\mathcal B}t_e\right)m^{-\tau}
   +\sum_{e\in\mathcal B}\left(\frac{t_e}{m}\right)^\tau
   \right].
\end{equation}
At $\tau=1$ this is exactly $s/(Wm)$: batching preserves the existing
rank deficit. The strict inequality $\mathcal C(\tau)<1$ is the
new time-exponent comparison. It uses only the original rank sum and
the count of selected ranks; no complete edge-rank histogram is needed.
Non-power widths use equal fields of width $\lfloor e/m\rfloor$,
at most $m-1$ separately handled high digits, and the same fixed-tape
row recursion. With sufficient rows the nonrecursive work is $O(V)$,
so the usual strong induction gives $O(Ve^\tau)$ whenever
\eqref{eq:batched-characteristic} is strict. The separate row and
arbitrary-width construction specifies this rounding step and removes
the initial row divisibility restriction.
